\section{Experimental Design}
\label{sec:sbc_experimental_design}

All SBC experiments were designed to represent the typical format of a \textit{Brasileirão Série A} season.
For each model, $S = 500$ synthetic datasets were generated, each comprising $n = 20$ teams competing in a double round-robin tournament with $380$ matches.

The inference process uses Stan \cite{carpenter_2017} via the CmdStanPy interface \cite{cmdstanpy}.
Each simulation was run with 4 parallel chains, 1,000 warm-up iterations, and 1,000 sampling iterations per chain, yielding 4,000 posterior samples per simulation.
The sampler was configured with an adaptation target acceptance probability of $\delta = 0.95$ and a maximum tree depth of 10.

For each simulation $i$, the true parameter values $\boldsymbol{\theta}_i$ were sampled from the prior distributions specified in the model.
Synthetic data $\mathbf{y}_i$ were generated according to the likelihood, and posterior samples $\{\boldsymbol{\theta}_i^{(1)}, \ldots, \boldsymbol{\theta}_i^{(L)}\}$ were obtained via Markov Chain Monte Carlo (MCMC).
The normalized rank of each true parameter value among its posterior samples was recorded.
Under correct posterior inference, these ranks should follow a uniform distribution on $[0, 1]$, as established in Theorem~\ref{thm:rank-uniformity}.

Since there are a lot of models, with many parameters, to evaluate, we use a systematic approach to identify potential miscalibrations.
For each parameter in a model, we compute the ECDF difference plot and check whether any point falls outside the 95\% confidence envelope.
If at least one point lies outside this envelope, that parameter is flagged, and we record the total number of points outside the envelope for that parameter.
Aggregating across all parameters of a model, we obtain two summary statistics: the total number of flagged parameters (``Flags'') and the cumulative count of points outside the confidence envelopes (``Points Out'').
A model with many flagged parameters or a large cumulative count suggests potential miscalibration.
However, under correct calibration, we expect approximately 5\% of parameters to be flagged by chance alone, and minimal deviations may still occur due to the inherent randomness of the simulations.
